% ---------------------------------------------------------------------
% VI. CONCLUSION AND LIMITATIONS (about 0.25 page)
% Owner: Rory Linnane
% ---------------------------------------------------------------------
\section{Conclusion and Limitations}\label{sec:conclusion}

We tested whether SHAP and LIME can support a social-bias audit of a pretraining data filter.
Our experiments show that both explainers are faithful enough to be useful: they identify words that genuinely drive removal decisions, with deletion of SHAP's top-3 tokens causing a mean log-odds drop of 3.09 versus 0.35 for random tokens, and their attributions align with counterfactual identity-swap effects and with exact Shapley values on a transparent filter.
However, faithfulness is not uniform: stability degrades on weak-signal sentences, attributions shift with the choice of removal baseline, and the two explainers disagree more often than their overall agreement rate implies~\cite{krishna2022disagreement}.
We therefore recommend that auditors use attribution as a triage tool --- to flag suspicious filter behaviour for further review --- rather than as a certification of fairness.
In practice: (1)~report LIME's surrogate $R^2$ and re-run with multiple seeds to check stability; (2)~treat high SHAP--LIME disagreement as a signal for human review; (3)~confirm any attribution finding with a counterfactual swap before drawing conclusions.

\smallskip
\noindent\textbf{Limitations.}
Our scope is purposely narrow: we use short English template sentences with single word identity substitutions, which do not capture dialect-level bias (e.g.\ AAE~\cite{sap2019risk}) where no single word carries the signal.
We test one filter family (Detoxify \texttt{original} and \texttt{unbiased}), and the two models differ in architecture (BERT vs.\ RoBERTa) as well as debiasing objective, so architectural effects cannot be fully separated.
The transparent filter in Experiment~4 is our own construction, not a real pipeline filter.
Our ground-truth references --- counterfactual effects and exact Shapley values --- exist only because we controlled the inputs; a real auditor facing an opaque commercial filter would not have them.
We did not evaluate quality filters, blocklists, or non-English data.

\smallskip
\noindent\textbf{Future work.}
Extending attribution audits to real web text from C4 or Common Crawl, to the quality filter of~\cite{gururangan2022whose}, and to human rationale benchmarks such as HateXplain~\cite{ngueajio2025decoding} would address the most significant scope limitations.
